\documentclass{beamer}

\usetheme{AnnArbor}
\usecolortheme{beaver}

\usepackage[utf8]{inputenc}
\usepackage[spanish]{babel}
\usepackage{amsmath}
\usepackage{tikz}
\usepackage{listings}
\usetikzlibrary{trees,positioning,arrows.meta}

\title{Fundamentos de Generación de Código}
\subtitle{Compiladores}
\author{}
\date{}

\lstset{
    basicstyle=\ttfamily\scriptsize,
    keywordstyle=\color{blue},
    commentstyle=\color{green!50!black},
    frame=single,
    breaklines=true,
    showstringspaces=false
}

\begin{document}

%------------------------------------------------
\begin{frame}
\titlepage
\end{frame}

%------------------------------------------------
\begin{frame}{Objetivo de la Clase}

Al finalizar la clase, el estudiante comprenderá:

\vspace{0.3cm}

\begin{itemize}
    \item el pipeline general de un compilador;
    \item qué hace la fase de generación de código;
    \item qué son AST, tripletas y cuádruplos;
    \item cómo se genera ensamblador;
    \item cómo se traducen expresiones y llamadas a función.
\end{itemize}

\end{frame}

%------------------------------------------------
\begin{frame}{Pipeline General del Compilador}

\centering

\begin{tikzpicture}[
node distance=1.0cm,
>=Stealth,
every node/.style={
draw,
rounded corners,
font=\scriptsize,
align=center,
minimum width=4.2cm
}
]

\node[fill=blue!15] (src) {Código Fuente\\programa.c};

\node[fill=green!20, below of=src] (lex) {
Análisis Léxico\\
Tokens
};

\node[fill=green!20, below of=lex] (syn) {
Análisis Sintáctico\\
AST / estructura
};

\node[fill=green!20, below of=syn] (sem) {
Análisis Semántico\\
Tipos, ámbitos, funciones
};

\node[fill=yellow!25, below of=sem] (ir) {
Representación Intermedia\\
Tripletas / Cuádruplos
};

\node[fill=orange!25, below of=ir] (gen) {
Generación de Código\\
Ensamblador
};

\node[fill=red!20, below of=gen] (exe) {
Código Objeto\\
Ejecutable
};

\draw[->] (src) -- (lex);
\draw[->] (lex) -- (syn);
\draw[->] (syn) -- (sem);
\draw[->] (sem) -- (ir);
\draw[->] (ir) -- (gen);
\draw[->] (gen) -- (exe);

\end{tikzpicture}

\end{frame}

%------------------------------------------------
\begin{frame}{¿Qué es Generación de Código?}

\begin{block}{Definición}
La generación de código transforma estructuras abstractas validadas por el compilador en instrucciones cercanas a la máquina.
\end{block}

\vspace{0.4cm}

La entrada del generador puede ser:

\begin{itemize}
    \item AST,
    \item tripletas,
    \item cuádruplos,
    \item código de tres direcciones.
\end{itemize}

\vspace{0.4cm}

La salida puede ser:

\begin{itemize}
    \item ensamblador,
    \item código objeto,
    \item ejecutables.
\end{itemize}

\end{frame}

%------------------------------------------------
\begin{frame}[fragile]{Ejemplo Base}

Usaremos:

\begin{lstlisting}[language=C]
x = a + b * c;
\end{lstlisting}

\vspace{0.4cm}

Este ejemplo permite estudiar:

\begin{itemize}
    \item precedencia;
    \item AST;
    \item tripletas;
    \item cuádruplos;
    \item generación de ensamblador.
\end{itemize}

\end{frame}

%------------------------------------------------
\begin{frame}{AST de la Expresión}

\centering

\begin{tikzpicture}[
level distance=1.1cm,
sibling distance=2.4cm,
every node/.style={
draw,
rounded corners,
fill=blue!12,
font=\small
}
]

\node {=}
child { node {x} }
child {
    node {+}
    child { node {a} }
    child {
        node {*}
        child { node {b} }
        child { node {c} }
    }
};

\end{tikzpicture}

\vspace{0.5cm}

\begin{block}{Interpretación}
Primero se calcula \texttt{b * c}, después se suma con \texttt{a}.
\end{block}

\end{frame}

%------------------------------------------------
\begin{frame}[fragile]{Código de Tres Direcciones}

Para:

\begin{lstlisting}[language=C]
x = a + b * c;
\end{lstlisting}

se genera:

\begin{lstlisting}
t1 = b * c
t2 = a + t1
x  = t2
\end{lstlisting}

\vspace{0.4cm}

Cada instrucción contiene una sola operación principal.

\end{frame}

%------------------------------------------------
\begin{frame}{Tripletas}

Forma general:

\[
(op,\; arg_1,\; arg_2)
\]

\vspace{0.3cm}

Para:

\[
x = a + b * c
\]

\centering

\begin{tabular}{|c|c|c|c|}
\hline
Índice & Operador & Arg1 & Arg2 \\
\hline
0 & * & b & c \\
\hline
1 & + & a & (0) \\
\hline
2 & = & x & (1) \\
\hline
\end{tabular}

\vspace{0.4cm}

\begin{block}{Idea}
Las tripletas referencian resultados previos usando índices.
\end{block}

\end{frame}

%------------------------------------------------
\begin{frame}{Cuádruplos}

Forma general:

\[
(op,\; arg_1,\; arg_2,\; resultado)
\]

\vspace{0.3cm}

\centering

\begin{tabular}{|c|c|c|c|c|}
\hline
Índice & Operador & Arg1 & Arg2 & Resultado \\
\hline
0 & * & b & c & t1 \\
\hline
1 & + & a & t1 & t2 \\
\hline
2 & = & t2 & -- & x \\
\hline
\end{tabular}

\vspace{0.4cm}

\begin{block}{Ventaja}
Los cuádruplos hacen explícitos los temporales.
\end{block}

\end{frame}

%------------------------------------------------
\begin{frame}[fragile]{Ejemplo: Tripleta a Ensamblador}

Tripleta:

\[
(*,\; b,\; c)
\]

\vspace{0.3cm}

Interpretación:

\[
resultado_0 = b * c
\]

\vspace{0.3cm}

Ensamblador conceptual:

\begin{lstlisting}
movl b, %eax
imull c, %eax
\end{lstlisting}

\vspace{0.3cm}

\begin{itemize}
    \item \texttt{b} se carga en \texttt{\%eax}.
    \item \texttt{c} multiplica el contenido del registro.
    \item El resultado queda en \texttt{\%eax}.
\end{itemize}

\end{frame}

%------------------------------------------------
\begin{frame}[fragile]{Ejemplo Completo: Tripletas a ASM}

Tripletas:

\begin{lstlisting}
0: (*, b, c)
1: (+, a, (0))
2: (=, x, (1))
\end{lstlisting}

\vspace{0.3cm}

Ensamblador conceptual:

\begin{lstlisting}
movl b, %eax
imull c, %eax

movl a, %ebx
addl %eax, %ebx

movl %ebx, x
\end{lstlisting}

\vspace{0.3cm}

La tripleta 0 produce el valor temporal utilizado por la tripleta 1.

\end{frame}

%------------------------------------------------
\begin{frame}[fragile]{Ejemplo: Cuádruplos a ASM}

Cuádruplos:

\begin{lstlisting}
(*, b, c, t1)
(+, a, t1, t2)
(=, t2, --, x)
\end{lstlisting}

\vspace{0.3cm}

Ensamblador conceptual:

\begin{lstlisting}
movl b, %eax
imull c, %eax
movl %eax, t1

movl a, %ebx
addl t1, %ebx
movl %ebx, t2

movl t2, x
\end{lstlisting}

\vspace{0.3cm}

\begin{block}{Observación}
Los cuádruplos generan ensamblador más explícito porque los temporales ya están definidos.
\end{block}

\end{frame}

%------------------------------------------------
\begin{frame}[fragile]{Ejemplo: Llamada a Función}

Código C:

\begin{lstlisting}[language=C]
r = suma(a, b);
\end{lstlisting}

Código intermedio:

\begin{lstlisting}
param a
param b
t1 = call suma, 2
r = t1
\end{lstlisting}

\end{frame}

%------------------------------------------------
\begin{frame}[fragile]{Llamada a Función a Ensamblador}

Ensamblador conceptual:

\begin{lstlisting}
movl a, %edi
movl b, %esi

call suma

movl %eax, r
\end{lstlisting}

\vspace{0.3cm}

\begin{itemize}
    \item Los argumentos se colocan en registros.
    \item \texttt{call} transfiere el control.
    \item El valor de retorno queda en \texttt{\%eax}.
\end{itemize}

\end{frame}

%------------------------------------------------
\begin{frame}{Bloques Básicos}

Un bloque básico es una secuencia de instrucciones con:

\begin{itemize}
    \item una sola entrada;
    \item una sola salida;
    \item sin saltos internos.
\end{itemize}

\vspace{0.4cm}

Ejemplo:

\[
t1 = b * c
\]

\[
t2 = a + t1
\]

\[
x = t2
\]

\end{frame}


%------------------------------------------------
\begin{frame}[fragile]{Demostración Terminal}

Generar ensamblador:

\begin{lstlisting}[language=bash]
gcc -S programa.c -o programa.s
\end{lstlisting}

Mostrar ensamblador:

\begin{lstlisting}[language=bash]
cat programa.s
\end{lstlisting}

Generar objeto:

\begin{lstlisting}[language=bash]
gcc -c programa.s -o programa.o
\end{lstlisting}

Generar ejecutable:

\begin{lstlisting}[language=bash]
gcc programa.o -o programa
\end{lstlisting}

\end{frame}

%------------------------------------------------
\begin{frame}{Resumen Conceptual}

\begin{itemize}
    \item El análisis léxico produce tokens.
    \item El análisis sintáctico produce estructura.
    \item El análisis semántico valida significado.
    \item La generación de código produce instrucciones ejecutables.
    \item Las tripletas usan índices.
    \item Los cuádruplos usan temporales explícitos.
    \item El ensamblador utiliza registros e instrucciones reales.
\end{itemize}

\end{frame}

%------------------------------------------------
\begin{frame}{Frase Final}

\begin{block}{Idea Central}
La generación de código transforma estructuras abstractas validadas por el compilador en instrucciones ejecutables para la arquitectura destino.
\end{block}

\vspace{0.5cm}

\centering

\Large

Código Fuente $\rightarrow$ AST $\rightarrow$ Tripletas / Cuádruplos $\rightarrow$ ASM $\rightarrow$ Ejecutable

\end{frame}

\end{document}
